% =========================================================================
%  Instrumentación Electrónica — PostLab de la sesión 1 (amplificadores operacionales)
%  Estilo heredado de los apuntes de Electromagnetismo I (Temas 0 y 1):
%  amsart, Palatino, versalitas, enunciados numerados por sección, paleta de
%  azules y grises sobrios. Ecuaciones numeradas a la derecha (opción reqno).
% =========================================================================

% --- IDIOMA Y CODIFICACIÓN ------------------------------------------------
\usepackage[utf8]{inputenc}
\usepackage[T1]{fontenc}
\usepackage[spanish,es-tabla,es-nodecimaldot]{babel}

% --- TIPOGRAFÍA (Palatino) --------------------------------------------------
\usepackage{mathpazo}
\linespread{1.05}
\emergencystretch=2em
\usepackage[final]{microtype}

% --- MATEMÁTICAS -------------------------------------------------------------
\usepackage{amsmath, amssymb, mathtools}
\usepackage{bm}

% --- UNIDADES (coma decimal, como en el resto de los apuntes) ---------------
\usepackage{siunitx}
\sisetup{output-decimal-marker={,}, group-separator={\,}, group-minimum-digits=5,
         exponent-product=\cdot, per-mode=symbol, range-phrase={--},
         list-final-separator={ y }, list-pair-separator={ y }}

% --- PÁGINA, COLORES Y GRÁFICOS --------------------------------------------
\usepackage[a4paper,margin=2.6cm]{geometry}
\usepackage{graphicx}
\usepackage{xcolor}
% Paleta base (idéntica a la de los apuntes de Electromagnetismo I)
\definecolor{azulprof}{HTML}{1F3A5F}    % azul oscuro
\definecolor{azulacero}{HTML}{3B6EA5}   % azul medio
\definecolor{azulclaro}{HTML}{8FB3D9}   % azul claro
\definecolor{azulfondo}{HTML}{F3F7FB}   % fondo tenue de definiciones
\definecolor{grispizarra}{HTML}{6B7280} % gris
\definecolor{grisclaro}{HTML}{D4D8DE}   % gris claro
\definecolor{acento}{HTML}{B8800A}      % acento: ámbar oscuro
\definecolor{rojoteo}{HTML}{B8800A}     % señales de salida, resultados
\definecolor{verdeteo}{HTML}{2F6F68}    % corrientes y sentidos de referencia
\definecolor{opfondo}{HTML}{E4EEF8}     % relleno del triángulo del AO

\usepackage{tikz}
\usetikzlibrary{babel, arrows.meta, calc, positioning, decorations.markings,
                patterns, intersections}
\usepackage[european, straightvoltages, RPvoltages]{circuitikz}
\usepackage{pgfplots}
\pgfplotsset{compat=1.17}
\usepgfplotslibrary{fillbetween}
\usepackage{booktabs}
\usepackage{needspace}
\usepackage{array}
\usepackage{multirow}
\usepackage[inline]{enumitem}
\setlist[itemize]{leftmargin=1.6em, itemsep=2pt, topsep=4pt}
\setlist[enumerate]{leftmargin=1.8em, itemsep=2pt, topsep=4pt}

% --- ENLACES Y REFERENCIAS CRUZADAS ----------------------------------------
\usepackage[colorlinks=true, linkcolor=azulacero, citecolor=azulacero,
            urlcolor=azulacero, plainpages=false, pdfusetitle=false,
            pdftitle={PostLab. Sesión 1: amplificadores operacionales},
            pdfauthor={Daniel Martín Guerrero}]{hyperref}
\usepackage[spanish,nameinlink]{cleveref}

% --- CABECERAS ---------------------------------------------------------------
\pagestyle{headings}
\makeatletter
\def\ps@headings{%
  \let\@mkboth\markboth
  \def\@evenhead{\small\scshape\normalfont\thepage\hfil Instrumentación Electrónica}%
  \def\@oddhead{\small\scshape\normalfont PostLab. Sesión 1: amplificadores operacionales\hfil\thepage}%
  \let\@evenfoot\@empty \let\@oddfoot\@empty}
\makeatother

% Subsecciones como títulos en bloque (no «run-in»), con aire antes y después
\usepackage{titlesec}
\titleformat{\subsection}[block]{\normalfont\bfseries\color{azulprof}}{\thesubsection.}{0.6em}{}
\titlespacing*{\subsection}{0pt}{3ex plus 1ex minus .5ex}{1.6ex plus .4ex}
\setcounter{tocdepth}{2}
\makeatletter
\def\@@and{y}
\renewcommand{\keywordsname}{Palabras clave}
% Índice: número de sección en una caja de anchura fija (evita el solape de amsart)
\renewcommand{\tocsection}[3]{\indentlabel{\@ifnotempty{#2}{\makebox[1.6em][l]{#2.}}}#3}
\renewcommand{\tocsubsection}[3]{\indentlabel{\@ifnotempty{#2}{\makebox[2.4em][l]{#2.}}}#3}
\def\l@section{\@tocline{1}{2pt}{0pt}{1.6em}{}}
\def\l@subsection{\@tocline{2}{0pt}{1.6em}{2.4em}{}}
\def\l@part{\@tocline{-1}{8pt plus 2pt}{0pt}{0pt}{\bfseries\color{azulprof}}}
\makeatother

% =========================================================================
%               NOTACIÓN Y MACROS DE USO FRECUENTE
% =========================================================================
\newcommand{\dif}{\mathrm{d}}
\newcommand{\eul}{\mathrm{e}}
\newcommand{\jj}{\mathrm{j}}                         % unidad imaginaria (j)
\newcommand{\R}{\mathbb{R}}
\newcommand{\C}{\mathbb{C}}
\newcommand{\abs}[1]{\lvert #1\rvert}
\newcommand{\defeq}{\mathrel{\overset{\mathrm{def}}{=}}}
\newcommand{\Lt}{\mathcal{L}}                        % transformada de Laplace
\DeclareMathOperator{\sen}{sen}
\DeclareMathOperator{\sgn}{sgn}
\DeclareMathOperator{\arctg}{arctg}
\DeclareMathOperator{\Real}{Re}
\renewcommand{\sin}{\sen}                            % «sen» en todo el texto
% Magnitudes de circuito
\newcommand{\vin}{v_{\mathrm{in}}}
\newcommand{\vout}{v_{\mathrm{out}}}
\newcommand{\Vin}{V_{\mathrm{in}}}
\newcommand{\Vout}{V_{\mathrm{out}}}
\newcommand{\Vsat}{V_{\mathrm{sat}}}
\newcommand{\Vcc}{V_{\mathrm{CC}}}
\newcommand{\Vee}{V_{\mathrm{EE}}}
\newcommand{\Rin}{R_{\mathrm{in}}}
\newcommand{\Rout}{R_{\mathrm{out}}}
\newcommand{\Zin}{Z_{\mathrm{in}}}
\newcommand{\Zout}{Z_{\mathrm{out}}}
\newcommand{\Av}{A_{V}}
\newcommand{\vp}{v_{+}}
\newcommand{\vm}{v_{-}}
\newcommand{\dB}{\,\mathrm{dB}}
\newcommand{\pend}[1]{\textcolor{grispizarra}{\colorbox{grisclaro}{\strut\,#1\,}}}

% =========================================================================
%        ENTORNOS FORMALES (enunciados numerados, estilo artículo)
% =========================================================================
\usepackage{amsthm}
\usepackage[most]{tcolorbox}

\newtheoremstyle{cienplano}{10pt}{10pt}{\itshape}{}%
  {\bfseries\color{azulprof}}{.}{0.6em}%
  {\thmname{#1}\thmnumber{ #2}\thmnote{\normalfont\upshape\ (#3)}}
\newtheoremstyle{ciendef}{10pt}{10pt}{\normalfont}{}%
  {\bfseries\color{azulprof}}{.}{0.6em}%
  {\thmname{#1}\thmnumber{ #2}\thmnote{\normalfont\upshape\ (#3)}}
\newtheoremstyle{cienobs}{8pt}{8pt}{\normalfont}{}%
  {\itshape\color{azulprof}}{.}{0.6em}%
  {\thmname{#1}\thmnumber{ #2}\thmnote{\normalfont\upshape\ (#3)}}

\theoremstyle{cienplano}
\newtheorem{proposicion}{Proposición}
\newtheorem{teorema}[proposicion]{Teorema}
\newtheorem{corolario}[proposicion]{Corolario}
\theoremstyle{ciendef}
\newtheorem{ejemplo}[proposicion]{Ejemplo}
\newtheorem{ejercicio}[proposicion]{PreLab}
\newtheorem{hipotesis}[proposicion]{Hipótesis}
% Definiciones: único entorno enmarcado (caja fina, fondo tenue, título en el borde)
\newcounter{definicion}
\renewcommand{\thedefinicion}{\theproposicion}
\newtcolorbox{definicionbox}[1]{%
  enhanced, breakable, sharp corners,
  colback=azulfondo, colframe=azulclaro, boxrule=0.4pt,
  left=9pt, right=9pt, top=5pt, bottom=5pt,
  before skip=12pt, after skip=12pt,
  attach boxed title to top left={yshift*=-\tcboxedtitleheight/2, xshift=8pt},
  boxed title style={colback=white, colframe=white, sharp corners, boxrule=0pt,
                     left=3pt, right=3pt, top=1pt, bottom=1pt},
  fonttitle=\bfseries\color{azulprof}, coltitle=azulprof,
  title={Definición~\thedefinicion\ifx\relax#1\relax\else\ (#1)\fi}}
\newenvironment{definicion}[1][]
  {\par\stepcounter{proposicion}\refstepcounter{definicion}\begin{definicionbox}{#1}}
  {\end{definicionbox}}

% Enunciado de cada recuadro del PreLab: misma caja, con filete ámbar (los
% recuadros amarillos del guion), para distinguir lo pedido de lo resuelto.
\newcounter{prelab}
\newtcolorbox{prelabbox}[1]{%
  enhanced, sharp corners,
  colback=white, colframe=acento!70, boxrule=0.5pt, leftrule=2.2pt,
  left=9pt, right=9pt, top=5pt, bottom=5pt,
  before skip=12pt, after skip=10pt,
  attach boxed title to top left={yshift*=-\tcboxedtitleheight/2, xshift=8pt},
  boxed title style={colback=white, colframe=white, sharp corners, boxrule=0pt,
                     left=3pt, right=3pt, top=1pt, bottom=1pt},
  fonttitle=\bfseries\color{acento!85!black}, coltitle=acento!85!black,
  title={PreLab~\theprelab\ifx\relax#1\relax\else\ (#1)\fi}}
\newenvironment{prelab}[1][]
  {\par\refstepcounter{prelab}\begin{prelabbox}{#1}}
  {\end{prelabbox}}

\theoremstyle{cienobs}
\newtheorem{observacion}[proposicion]{Observación}

% Resolución: misma maquinaria que «proof», con otro rótulo
\renewcommand{\qedsymbol}{\ensuremath{\blacksquare}}
\newenvironment{resolucion}[1][Resolución]
  {\begin{proof}[\itshape\color{azulprof}#1]}{\end{proof}}

% Resultado final enmarcado (como \boxed, pero con el filete de la paleta)
\newcommand{\resultado}[1]{\fcolorbox{azulacero}{white}{\ensuremath{\displaystyle #1}}}

\crefname{proposicion}{proposición}{proposiciones}
\Crefname{proposicion}{Proposición}{Proposiciones}
\crefname{teorema}{teorema}{teoremas}
\Crefname{teorema}{Teorema}{Teoremas}
\crefname{corolario}{corolario}{corolarios}
\Crefname{corolario}{Corolario}{Corolarios}
\crefname{definicion}{definición}{definiciones}
\Crefname{definicion}{Definición}{Definiciones}
\crefname{ejemplo}{ejemplo}{ejemplos}
\Crefname{ejemplo}{Ejemplo}{Ejemplos}
\crefname{hipotesis}{hipótesis}{hipótesis}
\Crefname{hipotesis}{Hipótesis}{Hipótesis}
\crefname{prelab}{PreLab}{PreLabs}
\Crefname{prelab}{PreLab}{PreLabs}
\crefname{observacion}{observación}{observaciones}
\Crefname{observacion}{Observación}{Observaciones}
\crefname{section}{sección}{secciones}
\Crefname{section}{Sección}{Secciones}
\crefname{subsection}{sección}{secciones}
\Crefname{subsection}{Sección}{Secciones}
\crefname{table}{tabla}{tablas}
\Crefname{table}{Tabla}{Tablas}
\crefname{figure}{figura}{figuras}
\Crefname{figure}{Figura}{Figuras}
\crefname{equation}{ecuación}{ecuaciones}
\Crefname{equation}{Ecuación}{Ecuaciones}
\crefformat{equation}{ecuación~(#2#1#3)}
\Crefformat{equation}{Ecuación~(#2#1#3)}
\crefrangeformat{equation}{ecuaciones~(#3#1#4)--(#5#2#6)}
\crefmultiformat{equation}{ecuaciones~(#2#1#3)}{ y~(#2#1#3)}{, (#2#1#3)}{ y~(#2#1#3)}

\usepackage[font=small, labelfont={sc,color=azulprof}, labelsep=period,
            margin=0.6cm, justification=justified]{caption}
\usepackage{subcaption}

% Referencia a un panel de figura: «figura 5\,(b)», enlazada en bloque
\newcommand{\figpanel}[2]{\hyperref[#1]{figura~\ref*{#1}\,\textup{(#2)}}}
\newcommand{\Figpanel}[2]{\hyperref[#1]{Figura~\ref*{#1}\,\textup{(#2)}}}

% =========================================================================
%        ESTILOS GRÁFICOS COMUNES (TikZ / circuitikz / pgfplots)
% =========================================================================
% =========================================================================
%  Estilos gráficos comunes: TikZ, circuitikz y pgfplots
% =========================================================================

% --- Estilos TikZ generales ------------------------------------------------
\tikzset{
  >=Stealth,
  eje/.style={->, line width=0.6pt, grispizarra},
  fino/.style={line width=0.4pt, grispizarra},
  punteado/.style={densely dotted, line width=0.5pt, grispizarra},
  etiq/.style={font=\small, text=black!85},
  nota/.style={font=\scriptsize, text=grispizarra, align=center},
  % corriente dibujada junto a un hilo (flecha verde con rótulo)
  corr/.style={-{Stealth[length=4.5pt,width=4pt]}, line width=0.8pt, verdeteo},
  % nudo con tensión anotada
  tens/.style={font=\small, text=azulprof},
}

% --- circuitikz: trazo y colores sobrios -----------------------------------
\ctikzset{
  bipoles/thickness=1,
  bipoles/length=1.15cm,
  tripoles/op amp/width=1.5,
  tripoles/op amp/height=1.35,
  tripoles/op amp/font=\small,
  current/distance=0.45,
  label/align=smart,
}
\tikzset{
  circuito/.style={line width=0.55pt, color=black!85, font=\small},
  ao/.style={fill=opfondo},
}

% --- Estilo base de ejes pgfplots ------------------------------------------
\pgfplotsset{
  tema/.style={
    tick label style={font=\scriptsize}, label style={font=\small},
    axis lines=left, axis line style={grispizarra},
    grid=major, grid style={grisclaro, line width=0.4pt},
    every tick/.style={grispizarra},
    legend style={font=\scriptsize, draw=none, fill=white, fill opacity=0.85,
                  text opacity=1},
    legend cell align=left,
    /pgf/number format/use comma,
  },
  % ejes centrados (formas de onda y características de transferencia)
  temacentro/.style={
    tick label style={font=\scriptsize}, label style={font=\small},
    axis lines=middle, axis line style={grispizarra, -{Stealth[length=4pt]}},
    every tick/.style={grispizarra},
    legend style={font=\scriptsize, draw=none, fill=white, fill opacity=0.85,
                  text opacity=1},
    legend cell align=left,
    /pgf/number format/use comma,
  },
  senal entrada/.style={azulacero, line width=1.0pt},
  senal salida/.style={rojoteo, line width=1.1pt},
  asintota/.style={grispizarra, densely dashed, line width=0.6pt},
}


% =========================================================================
%  Entornos y macros propios del PostLab
% =========================================================================
% Pregunta del guion (recuadros azules): caja con filete azul
\newcounter{pregunta}
\newtcolorbox{preguntabox}[1]{%
  enhanced, sharp corners,
  colback=azulfondo, colframe=azulclaro, boxrule=0.5pt, leftrule=2.2pt,
  left=9pt, right=9pt, top=5pt, bottom=5pt,
  before skip=12pt, after skip=10pt,
  attach boxed title to top left={yshift*=-\tcboxedtitleheight/2, xshift=8pt},
  boxed title style={colback=white, colframe=white, sharp corners, boxrule=0pt,
                     left=3pt, right=3pt, top=1pt, bottom=1pt},
  fonttitle=\bfseries\color{azulprof}, coltitle=azulprof,
  title={Cuestión~\thepregunta\ifx\relax#1\relax\else\ (#1)\fi}}
\newenvironment{pregunta}[1][]
  {\par\refstepcounter{pregunta}\begin{preguntabox}{#1}}
  {\end{preguntabox}}
\newenvironment{respuesta}[1][Respuesta]
  {\begin{proof}[\itshape\color{azulprof}#1]}{\end{proof}}
% Hueco para un oscilograma o gráfica con datos medidos
\newcommand{\hueco}[2]{%
  \tikz\node[draw=azulclaro, densely dashed, line width=0.6pt, fill=azulfondo,
             minimum width=#1, minimum height=3.6cm, align=center,
             font=\small, text=grispizarra, inner sep=4pt]{#2};}
% Dato pendiente de medir
\newcommand{\dato}{\textcolor{grisclaro}{\rule[-0.35ex]{1.3cm}{2.2ex}}}
